We present a hybrid quantum--classical framework for contextual sensor-data representation and evaluate its trainable quantum kernel on a controlled magnetometer task. The framework separates observations, eight-dimensional feature profiles, quantum encoding, and classical prediction. Its network demonstrator uses a regularized landmark embedding and a multilayer perceptron. The statistical benchmark instead isolates an eight-qubit fidelity kernel with 16 trainable parameters and a support vector classifier. A repository-preregistered protocol compares this model with radial-basis-function, neural, random-feature, and ensemble baselines across 30 independently generated corpora. Each corpus contains 120 episodes, partitioned into 72 training, 24 validation, and 24 test examples. Quantum natural gradient optimizes centered kernel--target alignment; validation selects the checkpoint and classifier regularization before test evaluation. Mean test balanced accuracy was 0.9528 for the quantum model and 0.9569 for the radial-basis-function baseline. The paired difference was -0.0042, with a 95\% bootstrap interval of [-0.0222, 0.0125] and a two-sided Wilcoxon p-value of 0.9328. The primary predictive-advantage criterion was not met. Optimization changed kernel spectra, but exploratory spectral--performance associations remained unresolved. Separate positive and training-label controls comprised 12 runs each. The latter retained validation supervision and failed the chance-compatibility check for both kernel methods. All quantum computations used exact classical simulation. The results support a traceable component-level evaluation, not a quantum hardware advantage or validation of the complete sensor-network demonstrator.
